\chapter{Entwicklung}
\label{c:dev}
Nach dem Abschluss der Konzeption kann zur Entwicklung über gegangen werden.
Die Rahmenbedingungen hierbei sind, dass der Hetzner Server, welcher in 
\ref*{s:konz:ser} beschrieben wurde verwendet wird und auf diesem ist Ubuntu
22 installiert.

\section{Vorbereitung}
Bevor mit der Entwicklung begonnen werden kann, müssen ein paar 
Dinge vorbereitet werden, damit man entiwckeln kann. 

\subsection{Modelle}
Der erste Schritt ist es die Modelle vorzubereiten.\par
Damit man die Modelle vorbereiten kann, wird als erstes 
Ollama installiert.
Dies kann über das Install Skript, welches auf Github zu finden \cite{OllamaOllama2024}
ist getan werden. Dazu werden noch die Nvidia Cuda Driver installiert, 
damit Ollama diese nutzen kann. Nun können die Modelle vorbereitet werden.\par
Hierfür werden die Modelle, die in \ref{s:konz:mod} vorgestellt
wurden über Huggingface gedownloadet.
Dies ist relativ simple über Git LFS zu erledigen.
Die Repositories der Modelle werden einfach über git clone kopiert.
Wenn die Dateien fertig gedownloadet sind, kann man die Modelle 
für das Importieren in Ollama vorbereiten.
Hierfür \par
Nachdem Ragas nicht mit den von Huggingface in Ollama geladenen 
Modellen funktioniert hat, wurde auf Ollama Modelle gewechselt. 
Die über \verb|Ollama pull model| geladen, wobei model der Name
des Models ist, welches man herunterladen möchte. Hierbei kann
man mittels \verb|model:latest| noch die Version angeben, welche man 
haben möchte. Hierbei kann man teils ebenfalls Größen angeben. Es kommt 
auf das jeweilige Model an, welche Optionen es gibt. Auf Ollama.com
findet man die jeweiligen Modelle und die Versionen, die man herunterladen kann.

\subsection{Google Cloud}
Hier geht es um die Vorbereitung des Google Cloud Projekts für die Verbindung 
mit der Google Chat API.\par
Der erste Schritt hierbei ist es, ein Google Cloud Projekt zu erstellen.
Wenn dies erledigt ist, kann man die Google Chat API aktivieren.\par
Dazu wird ein Google Service Account benötigt.
Die Konfiguration der API und des Service Accounts werden stattfinden, nachdem der Server konfiguriert wurde 
und das System funktionsfähig ist.

\subsection{Python Umgebung}

\subsection{Datenbank}
Hierfür wird auf dem Server Docker installiert.
Dazu wird \verb|chromadb/chroma| über \verb|docker pull| 
heruntergeladen und dann mittels \verb|docker run| gestartet.

\section{System}
Wie in der Konzeption beschrieben wird 
das System auf dem der Chatbot erstellt wird ist ein
Hetzner Server, auf dem ebenfalls die Datenbank liegt.
Nachdem die Vorbereitungen des Servers abgeschlossen sind
kann mit der eigentlichen Entwicklung begonnen werden.\par
\par





